\section{Task, Data and Protocol}
\label{sec:setup}

\paragraph{Prices.}
We use the CRSP daily stock file from January 1995 to December 2025.
The training universe defined below spans 7,550 trading days and 4,011 distinct stocks.
We keep ordinary common shares of U.S.\ companies listed on a U.S.\ exchange, using the share codes known at each date.
In the current CRSP format, a delisting return sits on a separate row that the standard common-stock filter removes; we keep these rows and set a missing delisting return to $-30\%$~\cite{shumway1997delisting}.
Nasdaq volumes are adjusted for double counting~\cite{gao2010marketing}.
Factor returns~\cite{fama2015five,carhart1997persistence} and the S\&P~500 are used only for reporting.

\paragraph{Accounting and analyst data.}
Quarterly accounting data come from the point-in-time Compustat tables, because later revisions replace the figures investors saw at the time.
Even there, the main value column matches the first-reported figure in only 60--70\% of quarters before 2015.
Taking the as-reported column when it is filled, and the main column otherwise, matches the first monthly snapshot in over 99.7\% of quarters from 1998 to 2024.
A quarter becomes usable only after that snapshot, and a late filing delays the quarters after it.
Analyst data use the \emph{unadjusted} I/B/E/S consensus for the current fiscal year, divided by the price on the date of the consensus, because the split-adjusted table rewrites past estimates for later splits~\cite{payne2003ibes}.
A consensus is usable from the day after its statistics date and an earnings surprise from the day after its announcement.
Records expire after 270, 45 and 100 days respectively.

\paragraph{Universes, target and timing.}
Each month-end, using only data known then, we build two universes that apply from the next session.
The \emph{training universe} holds the 1,000 largest stocks by market value with a price of at least \$5 and 252 days of history.
The \emph{tradable universe} holds the 100 stocks with the largest 63-day median dollar volume among the 500 largest.
A signal computed after the close of day $t$ is executed at the close of $t{+}1$, and the target is the return from the close of $t{+}1$ to the close of $t{+}6$, minus its mean over the training universe $U_t$:
\begin{equation}
y_{i,t} = R_{i,t} - \frac{1}{|U_t|}\sum_{j\in U_t} R_{j,t},\qquad R_{i,t} = \prod_{k=2}^{6} (1 + r_{i,t+k}) - 1.
\label{eq:target}
\end{equation}
Models learn the percentile rank of $y_{i,t}$ within each date.

\paragraph{Portfolio and costs.}
The long-only portfolio targets 20 equal-weight stocks, with decisions taken every day.
A holding is sold only when its rank in the tradable universe falls below 40, and the best-ranked stocks replace it.
A position above 7.5\% is cut back to 5\%, and a delisted holding moves to cash.
Every trade pays 2, 5 or 10 basis points (bp) per side, 5\,bp being the central scenario; a \emph{measured} scenario charges each stock's closing quoted half-spread plus 1\,bp.
The benchmark is the equal-weight tradable universe, which earned 10.7\% per year over 2005--2021, like the S\&P~500.

\begin{figure}[t]
\centering
% Evaluation protocol: data periods, one walk-forward fold, and the timing of one decision.
\begin{tikzpicture}[x=0.25cm, y=1cm, font=\footnotesize, every node/.style={inner sep=1pt}]
  % Data periods (x = years since 1995).
  \fill[black!10] (0,0) rectangle (10,0.45);
  \fill[cDev] (10,0) rectangle (27,0.45);
  \fill[cFinal] (27,0) rectangle (31,0.45);
  \node at (5,0.225) {warm-up, first training};
  \node at (18.5,0.225) {development: 17 yearly folds};
  \node[anchor=south] at (29,0.47) {final test, once};
  \foreach \x/\label in {0/1995, 10/2005, 27/2022, 31/2026}
    \node[anchor=north, text=black!65] at (\x,-0.02) {\label};

  % One fold: test year 2010, with purges between phases.
  \begin{scope}[yshift=-1.05cm]
    \fill[black!10] (0,0) rectangle (11.85,0.38);
    \fill[black!22] (12,0) rectangle (14.85,0.38);
    \fill[cDev!100!black!35] (15,0) rectangle (16,0.38);
    \node at (5.9,0.19) {train candidates};
    \node[anchor=south] at (13.4,0.4) {validation};
    \node[anchor=west] at (16.2,0.19) {test year 2010};
    \fill[black!10] (0,-0.5) rectangle (14.85,-0.12);
    \node at (7.4,-0.31) {refit on all data before 2010};
    \draw[black!55, -{Stealth[length=3pt]}] (16.9,-0.45) node[anchor=west, text=black!65] {6-session purges} -- (15.05,-0.31);
  \end{scope}

  % Timing of one decision, on its own scale.
  \begin{scope}[yshift=-2.5cm, x=1.05cm]
    \draw[black!45] (0,0) -- (6.4,0);
    \foreach \k in {0,...,6} \draw[black!45] (\k,-0.05) -- (\k,0.05);
    \node[anchor=north, align=center] at (0,-0.07) {$t$\\signal};
    \node[anchor=north, align=center] at (1,-0.07) {$t{+}1$\\trade};
    \node[anchor=north, align=center] at (6,-0.07) {$t{+}6$\\target ends};
    \draw[cCand, line width=1pt, {Stealth[length=3pt]}-{Stealth[length=3pt]}] (1,0.17) -- (6,0.17);
    \node[anchor=south] at (3.5,0.19) {target: return from the close of $t{+}1$ to the close of $t{+}6$};
  \end{scope}
\end{tikzpicture}

\caption{Evaluation protocol. Top: data periods. Middle: one development fold (test year 2010); hyperparameters are chosen on the three preceding years before a refit on all earlier data, with a 6-session purge between phases. Bottom: timing of one decision.}
\Description{A timeline from 1995 to 2026 with three periods: warm-up and first training from 1995 to 2004, development with 17 yearly folds from 2005 to 2021, and a final test run once from 2022 to 2025. Below, one fold shows training data, three validation years and the test year 2010, separated by purges. At the bottom, a decision timeline shows the signal at day t, the trade at t+1 and the target return from t+1 to t+6.}
\label{fig:protocol}
\end{figure}

\paragraph{Walk-forward protocol.}
Each development year from 2005 to 2021 (17 folds) is predicted by a model trained on all earlier data, with hyperparameters chosen on the three preceding years by validation \ric{} (\cref{fig:protocol}).
The final test, 2022--2025, is refused by the code unless explicitly unlocked, and was run once with a configuration recorded before any of its results were computed.

\paragraph{Metrics and selection rules.}
The \ric{} is the Spearman correlation between scores and targets across tradable stocks, averaged over dates, with Newey--West $t$-statistics (10 lags)~\cite{newey1987simple}.
Paired comparisons use a moving-block bootstrap~\cite{kunsch1989jackknife} of daily differences (21-day blocks, 2,000 draws).
For portfolios we report the annualized net active return, the information ratio (IR), the six-factor alpha and the one-way turnover.
All 70 development runs are registered, and the Deflated Sharpe ratio (DSR)~\cite{bailey2014deflated} uses their number and the variance of their Sharpe ratios.
Rules fixed beforehand: models are compared by net IR at 5\,bp; a more complex model replaces a simpler one only if a bootstrap interval of their \ric{} or return difference excludes zero; and a result is established only if its $t$ exceeds 2 \emph{and} its DSR exceeds 0.95.
